\section{背景知识}

\subsection{这个项目做什么}

TeXlate 是「幻觉翻译」(hjfy.top) 的开源复刻：输入一篇 arXiv 论文的 LaTeX 源码（e-print），输出一份重编译的中文 PDF，正文逐段替换为中文译文、数学公式/图表/参考文献结构原样保留，并提供双语对照阅读。链路上的每一段都由自研组件构成：

\begin{itemize}
  \item \textbf{获取层}：arXiv e-print 批量下载（IA bulk / export / OAI-PMH 三通道），钉版入库；
  \item \textbf{解析层}：自研半解析器——不是完整 TeX 引擎，而是「逐字符扫描 $\to$ pieces $\to$ 占位符保护 $\to$ 区间 splice」的保真切分器，外加宏展开机（Gullet）把用户自定义宏就地展开；
  \item \textbf{翻译层}：段落级 chunk 经 LLM 网关批量翻译（swe-2 系列模型），占位符协议保护公式与命令不被译文污染；
  \item \textbf{校验层}：L0/L1/L2 三层校验（chunk 不变量、tree-sitter 结构、LLM judge）；
  \item \textbf{编译层}：ctex 注入 + normalize + xelatex/tectonic 双引擎编译；
  \item \textbf{修复层}：fixloop——131 条 YAML 规则驱动的自动修复循环（装包、改源、路由、polyfill），把编译失败的论文迭代救回；
  \item \textbf{对齐层}：named-destination 锚点同步，滚动阅读时中英逐段对齐。
\end{itemize}

\subsection{为什么这件事难}

\textbf{arXiv 源码是脏的真实世界数据}。语料横跨 LaTeX2.09（1990s）到 2026 年：混用 plain \TeX{} 原语与 LaTeX2.09 旧式、非 UTF-8 编码（Big5/GBK/Latin-1）、缺文件的残缺包、tar/gz 双格式、古早宏包（revtex4、aastex、psfig）。基线实测：裸编译（无修复）union pdf 仅 65--72\%，clean 22--40\%——三分之二的论文源码直接编译不出 PDF，最大失败类是 missing\_file（占裸失败 92.1\%）。

\textbf{宏展开是现成解析器的共同盲区}。95\% 的论文含自定义宏（中位 47 个/篇，最多 446），其中 12/39 篇手挑语料存在「宏体里藏 \verb|\begin/\end|」的结构性陷阱，49\% 的宏体藏 $>20$ 字母的可译文本。对 8 个第三方解析库（latex-utensils、unified-latex、tree-sitter-latex、TexSoup、plasTeX、pylatexenc、LaTeX.js、ieeA）的横评（\texttt{bench/results/00-grand-comparison.md}）表明：陷阱断言全军覆没——最好成绩 24/26 vs 自研半解析器 32/32；「能 parse」不等于「parse 对」，pylatexenc 报 90/90 成功却是静默截断零报错。\textbf{无错误信号比崩溃更危险}，这决定了主解析器必须自研、校验器必须独立。

\textbf{翻译约束是结构性的}。chunk 边界纪律要求公式、命令、环境标签逐字节保留，译文只能在散文区间替换；术语表一致性、学术语域（register）、不可译段（参考文献/版权行）都要显式协议。编译侧中文需要 ctex 注入与 CJK 字体路由，修复侧要处理 EPS 图源、缺 sty/cls、bib 链路断点等 131 类已建档机制。

\subsection{管线架构}

图~\ref{fig:pipeline} 给出全链架构与各段的 benchmark 对应关系。

\begin{figure}[htbp]
\centering
\resizebox{\textwidth}{!}{%
\begin{tikzpicture}[
  font=\small,
  stage/.style={draw,rounded corners=2pt,minimum height=8mm,inner xsep=4pt,align=center,fill=cA!8},
  impl/.style={font=\scriptsize\color{black!60},align=center},
  bench/.style={draw=cC,dashed,rounded corners=1pt,font=\scriptsize\color{cC},inner sep=2pt,align=center,fill=cC!6},
  arr/.style={-{Stealth[length=2mm]},thick,black!70},
  node distance=5mm and 4mm
]
% 主流水线
\node[stage] (fetch) {arXiv\\e-print};
\node[stage,right=of fetch] (parse) {flatten + 半解析\\mouth/gullet/segmenter};
\node[stage,right=of parse] (xlat) {段落级翻译\\xlat 编排 + 网关};
\node[stage,right=of xlat] (valid) {校验\\L0/L1/L2};
\node[stage,right=of valid] (splice) {splice 重建\\+ ctex 注入};
\node[stage,right=of splice] (comp) {编译\\xelatex / tectonic};
\node[stage,right=of comp] (align) {锚点对齐\\+ 双语 PDF};
\draw[arr] (fetch) -- (parse);
\draw[arr] (parse) -- (xlat);
\draw[arr] (xlat) -- (valid);
\draw[arr] (valid) -- (splice);
\draw[arr] (splice) -- (comp);
\draw[arr] (comp) -- (align);
% fixloop 回环
\node[stage,below=6mm of comp,fill=cB!10,draw=cB] (fix) {fixloop 修复循环\\131 条 YAML 规则};
\draw[arr,cB] (comp.south) -- (fix.north);
\draw[arr,cB] (fix.west) -| (splice.south);
% bench 标注
\node[bench,below=4mm of parse] (b1) {B1 parsebench\\identity/leak/flatten};
\node[bench,below=4mm of xlat] (b4) {B4 xlatbench + qualbench\\硬契约/延迟/ESA 评分};
\node[bench,below=4mm of valid] (b6) {B6 validbench\\检出率/FP};
\node[bench,below=21mm of splice] (b3) {B3 compilebench\\clean/pdf 率};
\node[bench,below=4mm of align] (b7) {B7 alignbench\\named-dest 保留};
\node[bench,above=4mm of xlat] (b5) {B5 e2ebench 全链漏斗};
\draw[cC,dashed,-{Stealth[length=1.5mm]}] (b5.south) -- (xlat.north);
\node[bench,above=4mm of parse] (b2) {B2 fixtures T01--T98};
\end{tikzpicture}}%
\caption{TeXlate 管线架构与 benchmark 臂对应关系。主流水线自左向右；fixloop 是编译段的自动修复回环；B1--B7 为七个评测臂（虚线框），另有 fixtures 陷阱断言集做解析段单元级测试。}
\label{fig:pipeline}
\end{figure}

\subsection{术语与口径}

\begin{itemize}
  \item \textbf{cell / 格}：语料库中一篇论文的存储单元（\texttt{\{id\}/\{meta.json, raw.*, extracted/\}}）；benchmark 语境下也指一个被评对象（paper $\times$ engine 组合）。
  \item \textbf{层 / layer}：corpus\_v3 的抽样分层（见 \S\ref{sec:corpus}），每层独立 manifest、独立抽样框。
  \item \textbf{臂 / arm}：同一被测对象的不同处理条件——compile 的 baseline（裸编译）vs zh（产品链 normalize 后）；e2e 的 mock（假翻译）vs real（真 LLM）；stagerun 的 base vs zh 条件。
  \item \textbf{clean / pdf\~ / FAIL}：编译终态分级——clean 无 warning 出 PDF；pdf\~ 带可接受瑕疵出 PDF；FAIL 无 PDF。union pdf = clean + pdf\~。
  \item \textbf{identity / leak}：解析质量两轴——identity 为 reconstruct 逐字节一致率；leak 为受保护内容（公式/命令）漏进可译 chunk 的占比。
  \item \textbf{scorecard}：gate\_scorecard.py 对 stagerun records 的终态统计，M2 出口门为 union pdf $\geq 90\%$、clean $\geq 90\%$。
  \item \textbf{机制 / mechanism}：mechanisms.jsonl 台账登记的一类已知缺陷模式（如 dollar 族泄漏、eps\_route），benchmark 发现 $\to$ 归因 $\to$ 规则沉淀的反馈环。
\end{itemize}
